% Author component. Root alone incorporates this into the receiving manuscript.
% The 198 arithmetic and Busy Beaver proofs are retained without alteration.
\section{Brox's elimination from 198 to 196 coordinates}\label{q196:elimination}
We apply the monomial elimination of Jos\'e Brox~\cite{q196:brox}
to the 198-coordinate system proved above. This is an equivalence of
entire rational solution sets for every fixed parameter; it changes
neither the arithmetic hypotheses nor the sufficient parameter.

\begin{lemma}[Elimination with nonzero rescaling coordinates]\label{q196:general}
Let a polynomial system over $\Q$ have input $x$, retained coordinates
$u_1,\ldots,u_k$, further coordinates $x_i$, and coordinates $t_j$
to be eliminated. Suppose retained equations imply
$u_1\cdots u_k\ne0$, and the equations defining $t_j$ uniquely express
them as Laurent monomials in the $u_\ell$ with nonzero rational
coefficients. Fix the input and the $u_\ell$, substitute those expressions,
and make the changes
\[
 x_i=c_i\prod_{\ell=1}^k u_\ell^{e_{i\ell}}z_i,
 \qquad c_i\in\Q^\times,\quad e_{i\ell}\in\Z.
\]
Suppose each remaining equation, after substitution and use of retained
power identities, is a nonzero rational constant times a Laurent
monomial in the $u_\ell$ times a new polynomial equation.
Suppose the new equations imply the same nonvanishing and power identities.
Then the old and new rational zero sets are in input-preserving bijection.
If the new equations have degree at most two, their primitive integer
forms have a sum of squares of degree at most four with the same input
projection.
\end{lemma}
\begin{proof}
At every old solution the nonzero coordinates permit inverse diagonal
scaling. The row identities give a solution of the new system.
At every new solution the required coordinates are again nonzero, so
diagonal scaling and the unique Laurent formulas recover all old
coordinates. The eliminated equations hold by construction, and each
remaining old equation is zero by its row identity. The two maps are
inverse and fix the input. Multiplying an equation by a nonzero rational
number does not change its zeros. Over $\Q$ a sum of squares vanishes
exactly when every summand vanishes.
\end{proof}

\paragraph{The three nonzero coordinates.}
For this paragraph use the parent's one-based indices
$x_i=Y_{i-1}$ from Section~\ref{q198:rows}. Put
\[
 R=x_{116},\quad J=x_{104},\quad u=x_{45},\quad v=x_{103},
 \qquad w_0=x_{78},\quad w_1=x_{93}.
\]
The following rows remain literally unchanged after renumbering:
\begin{align*}
 u&=x^2,\\
 4v&=6255001+255102u+2601u^2,\\
 0&=14315517724+291790992u+15531948uv+767410608v-512v^2+RJ,\\
 0&=1+22201RJ-AJ.
\end{align*}
Substituting the first two equations into the third and expanding gives
\[
 RJ=r(x):=(51x^2+103)(51x^2+2503)(51x^2+2504)(1632x^2+80471)>0.
\]
Thus $R,J\ne0$ for every rational input; for $A>0$ the last row gives
$R=Ar(x)/(1+22201r(x))>0$.
The remaining retained rows used for elimination are
\[
 x_{115}=R^2,\qquad x_{77}w_0=1,\quad x_{79}=w_0^2,
 \qquad x_{92}w_1=1,\quad x_{94}=w_1^2.
\]
In particular $w_0w_1\ne0$. These nine equations hold in the new system
as well, so there is no unexamined zero-denominator component.

\paragraph{The exact change of coordinates.}
Delete coordinates 80 and 95 and their two defining rows, using
\[
 x_{80}=w_0^4,\qquad x_{95}=R^{-1}w_1^4.
\]
Keep the remaining coordinates in their parent order. If $z_i$ denotes
the new coordinate carrying parent identifier $i$, set $x_i=m_i z_i$.
All unlisted multipliers are one. The nonconstant multipliers are
\[
\begin{array}{c|c@{\qquad}c|c}
 i&m_i&i&m_i\\\hline
 4,86&w_0^{-1}&8,101&w_1^{-1}\\
 5&w_0&9&w_1\\
 85&w_0^2&100&w_1^2\\
 81,82,83,84&w_0^{-2}&96,97,98,99&w_1^{-2}\\
 121,\ldots,136&w_0^{-2}&137,\ldots,146&Rw_1^{-2}\\
 &&147,148&w_1^{-2}\\
 &&149,150,151,152&R^{-1}w_1^{-2}
\end{array}
\]
For the constant multipliers put $H=31581660^{14}$ and $K=82479248H$.
They are
\[
\begin{array}{c|c}
 i&m_i\\\hline
 107&H^{-1}\\
 108&-K^{-1}\\
 11&-K\\
 12,13,14&K\\
 109,111&K^2\\
 110&-1
\end{array}
\]
These changes fix the input and the nine guard rows. None of $R,w_0,w_1$
is rescaled. The two removed defining equations become identities.
For each of the other 182 rows, substitute these expressions and reduce
only by the displayed retained power identities. Collecting equal
Laurent monomials gives a nonzero Laurent monomial in $R,w_0,w_1$
times a nonzero rational constant times the corresponding row printed
in Section~\ref{q196:rows}. Table~\ref{q196:rowmultipliers} gives the
source row, target row and multiplier for every such identity.
This is a coefficientwise identity: the
printed parent and new rows and the preceding multiplier tables specify
both sides completely. It can be checked by expanding each term,
collecting its powers of $R,w_0,w_1$ and its remaining coordinates,
then comparing the resulting rational coefficients. No arithmetic
equation beyond the nine displayed guard and power equations is used
in this calculation.

\begin{proposition}[The 196-coordinate family]\label{q196:family}
For every fixed integer $N\ge1$, put $A_N=149^{2N+1}$ and form
\[
 \mathcal P_{196,A_N}=\sum_{i=1}^{182}(q_i^{[196]}(A_N))^2
\]
from Section~\ref{q196:rows}. Its rational input projection equals that
of $\mathcal P_{198,A_N}$ from Section~\ref{q198:rows}.
It has 196 total coordinates, degree four, content one and 2943 nonzero
coordinate monomials. Its exact coefficient height is
\[
 H(\mathcal P_{196,A_N})=
 \max\{484991541504995159662206970923766531374664106904986910720,
       149^{4N+2}\}.
\]
\end{proposition}
\begin{proof}
The row identities and the nonvanishing just proved satisfy
Lemma~\ref{q196:general}. Exactly two coordinates and two rows were
deleted, giving $198-2=196$ coordinates and $184-2=182$ quadratics.
Expanding their sum of squares gives 2943 distinct coordinate monomials.
Every coordinate occurs, and the coefficient of the fourth power of the
input is one, proving both exact degree and integer content.
In the new table's zero-based coordinates $\mathbf Y$, the entire
parameter dependence is
\[
 \mathcal P_{196,A}(\mathbf Y)=\mathcal P_{196,0}(\mathbf Y)
 -2AY_{101}-44402AY_{101}^{2}Y_{113}+A^2Y_{101}^{2}.
\]
The three displayed monomials do not occur in $\mathcal P_{196,0}$.
The largest absolute coefficient in that constant part is the integer
in the asserted maximum. These finite assertions follow by multiplying
each printed row by itself and collecting equal monomials; the operation
uses integer arithmetic throughout. Since $A_N\ge149^3>44402$,
$A_N^2$ dominates the other parameter-dependent coefficients. No
coefficient cancels at positive $A_N$, proving the monomial count and
height formula for every $N\ge1$.
\end{proof}

The fixed integer in this maximum lies strictly between $149^{26}$
and $149^{27}$. Hence $N\ge7$ suffices for the coefficient height to be
$A_N^2$. This is only a coefficient crossover, not a defining threshold.

\paragraph{Transfer of the fixed Busy Beaver specialization.}
Theorem~\ref{q198bb:fixed} supplies the parent polynomial at
$N_{\rm BB}=\BB(2^{1000})$. Proposition~\ref{q196:family} applies
at that same parameter, so the new polynomial has exactly the same
rational input projection, namely $\Z$. The machine convention and
the threshold proof are unchanged. In particular $N_{\rm BB}\ge7$,
and the new coefficient height is the parent's large-parameter value
$A_{\rm BB}^2=149^{4\BB(2^{1000})+2}<\BB(2^{2000})$.
This proves Theorem~\ref{q196:fixed} without a new arithmetic threshold
or a new Busy Beaver program.
